\par\Needspace{6\baselineskip}\medskip
\begin{center}\hypertarget{native-181-article135}{}\pdfbookmark[2]{135.}{article-135}135.\end{center}
The preceding investigations were presented so that the principles on which our method rests
and, so to speak, its very core might be seen more clearly. The treatment itself, however, will
present the method in an entirely different form, which after very frequent use we can recommend
as the most convenient of the several forms we have tried. Since determining an unknown orbit
from three observations always comes down to a few successive hypotheses, or rather successive
approximations, it will be a considerable gain if we can arrange the calculation so that from
the outset we can separate from the hypotheses themselves as many operations as possible that
depend on a combination of known quantities alone, independently of $P$ and $Q$. These
preliminary operations, common to each hypothesis, then clearly need to be performed only once,
and the hypotheses themselves are reduced to as few operations as possible. It will be equally
important if we can avoid carrying each hypothesis as far as the elements themselves and reserve
their calculation for the \hypertarget{native-182}{}final hypothesis. In both respects our method, which we now begin
to explain, appears to leave nothing to be desired.
